\chapter{2017年全国II卷（文）}

\section{单选题}

\begin{problem}
设集合 \(A=\{1,2,3\}\)，\(B=\{2,3,4\}\)，则 \(A\cup B=\)（\quad）

\choices
    {\(\{1,2,3,4\}\)}
    {\(\{1,2,3\}\)}
    {\(\{2,3,4\}\)}
    {\(\{1,3,4\}\)}
\end{problem}

\begin{problem}
\((1+\i)(2+\i)=\)（\quad）

\choices
    {\(1-\i\)}
    {\(1+3\i\)}
    {\(3+\i\)}
    {\(3+3\i\)}
\end{problem}

\begin{problem}
函数 \(f(x)=\sin\!\left(2x+\frac{\pi}{3}\right)\) 的最小正周期为（\quad）

\choices
    {\(4\pi\)}
    {\(2\pi\)}
    {\(\pi\)}
    {\(\dfrac{\pi}{2}\)}
\end{problem}

\begin{problem}
设非零向量 \(\bs a\)，\(\bs b\)，满足 \(\left|\bs a+\bs b\right|=\left|\bs a-\bs b\right|\)，则（\quad）

\choices
    {\(\bs a\perp \bs b\)}
    {\(\left|\bs a\right|=\left|\bs b\right|\)}
    {\(\bs a\parallel\bs b\)}
    {\(\left|\bs a\right|>\left|\bs b\right|\)}
\end{problem}

\begin{problem}
若 \(a>1\)，则双曲线 \(\dfrac{x^2}{a^2}-y^2=1\) 的离心率的取值范围是（\quad）

\choices
    {\((\sqrt2,+\infty)\)}
    {\((\sqrt2,2)\)}
    {\((1,\sqrt2)\)}
    {\((1,2)\)}
\end{problem}

\begin{problem}
如图，网格纸上小正方形的边长为 \(1\)，粗实线画出的是某几何体的三视图，该几何体由一平面将一圆柱截去一部分后所得，则该几何体的体积为（\quad）

\bitmapfigure[float=inline,width=0.68\linewidth]{img/2017/national_paper_2_liberal/q6_fig1.png}

\choices
    {\(90\pi\)}
    {\(63\pi\)}
    {\(42\pi\)}
    {\(36\pi\)}
\end{problem}

\begin{problem}
设 \(x\)，\(y\) 满足约束条件
\(\begin{cases}
2x+3y-3\le 0,\\
2x-3y+3\ge 0,\\
y+3\ge 0
\end{cases}\)
则 \(z=2x+y\) 的最小值是（\quad）

\choices
    {\(-15\)}
    {\(-9\)}
    {\(1\)}
    {\(9\)}
\end{problem}

\begin{problem}
函数 \(f(x)=\ln(x^2-2x-8)\) 的单调递增区间是（\quad）

\choices
    {\((-\infty,-2)\)}
    {\((-\infty,-1)\)}
    {\((1,+\infty)\)}
    {\((4,+\infty)\)}
\end{problem}

\begin{problem}
甲，乙，丙，丁四位同学一起去向老师询问成语竞赛的成绩. 老师说：你们四人中有 \(2\) 位优秀，\(2\) 位良好，我现在给甲看乙，丙的成绩，给乙看丙的成绩，给丁看甲的成绩. 看后甲对大家说：我还是不知道我的成绩. 根据以上信息，则（\quad）

\choices
    {乙可以知道四人的成绩}
    {丁可以知道四人的成绩}
    {乙，丁可以知道对方的成绩}
    {乙，丁可以知道自己的成绩}
\end{problem}

\begin{problem}
执行如图的程序框图，如果输入的 \(a=-1\)，则输出的 \(S=\)（\quad）

\bitmapfigure[float=inline,width=0.68\linewidth]{img/2017/national_paper_2_liberal/q10_fig1.png}

\choices
    {\(2\)}
    {\(3\)}
    {\(4\)}
    {\(5\)}
\end{problem}

\begin{problem}
从分别写有 \(1\)，\(2\)，\(3\)，\(4\)，\(5\) 的 \(5\) 张卡片中随机抽取 \(1\) 张，放回后再随机抽取 \(1\) 张，则抽得的第一张卡片上的数大于第二张卡片上的数的概率为（\quad）

\choices
    {\(\dfrac{1}{10}\)}
    {\(\dfrac{1}{5}\)}
    {\(\dfrac{3}{10}\)}
    {\(\dfrac{2}{5}\)}
\end{problem}

\begin{problem}
过抛物线 \(C:y^2=4x\) 的焦点 \(F\)，且斜率为 \(\sqrt3\) 的直线交 \(C\) 于点 \(M\)（\(M\) 在 \(x\) 轴上方），\(l\) 为 \(C\) 的准线，点 \(N\) 在 \(l\) 上，且 \(MN\perp l\)，则 \(M\) 到直线 \(NF\) 的距离为（\quad）

\choices
    {\(\sqrt5\)}
    {\(2\sqrt2\)}
    {\(2\sqrt3\)}
    {\(3\sqrt5\)}
\end{problem}

\section{填空题}

\begin{problem}
函数 \(f(x)=2\cos x+\sin x\) 的最大值为\fillinblank{}.
\end{problem}

\begin{problem}
已知函数 \(f(x)\) 是定义在 \(\R\) 上的奇函数，当 \(x\in(-\infty,0)\) 时，\(f(x)=2x^3+x^2\)，则 \(f(2)=\fillinblank{}\).
\end{problem}

\begin{problem}
长方体的长，宽，高分别为 \(3\)，\(2\)，\(1\)，其顶点都在球 \(O\) 的球面上，则球 \(O\) 的表面积为\fillinblank{}.
\end{problem}

\begin{problem}
\(\triangle ABC\) 的内角 \(A\)，\(B\)，\(C\) 的对边分别为 \(a\)，\(b\)，\(c\)，若 \(2b\cos B=a\cos C+c\cos A\)，则 \(B=\fillinblank{}\).
\end{problem}

\section{解答题}

\begin{problem}
已知等差数列 \(\{a_n\}\) 的前 \(n\) 项和为 \(S_n\)，等比数列 \(\{b_n\}\) 的前 \(n\) 项和为 \(T_n\)，\(a_1=-1\)，\(b_1=1\)，\(a_2+b_2=2\).
\begin{enumerate}
\item 若 \(a_3+b_3=5\)，求 \(\{b_n\}\) 的通项公式；
\item 若 \(T_3=21\)，求 \(S_3\).
\end{enumerate}
\end{problem}

\begin{problem}
如图，四棱锥 \(P-ABCD\) 中，侧面 \(PAD\) 为等边三角形且垂直于底面 \(ABCD\)，\(AB=BC=\dfrac12AD\)，\(\angle BAD=\angle ABC=90^\circ\).
\begin{enumerate}
\item 证明：直线 \(BC\parallel\) 平面 \(PAD\)；
\item 若 \(\triangle PCD\) 面积为 \(2\sqrt7\)，求四棱锥 \(P-ABCD\) 的体积.
\end{enumerate}

\bitmapfigure[float=inline,width=8.5cm]{img/2017/national_paper_2_liberal/q18_fig1.png}
\end{problem}

\begin{problem}
海水养殖场进行某水产品的新，旧网箱养殖方法的产量对比，收获时各随机抽取了 \(100\) 个网箱，测量各箱水产品的产量（单位：\(\mathrm{kg}\)），其频率分布直方图如下：
\begin{enumerate}
\item 记 \(A\) 表示事件“旧养殖法的箱产量低于 \(50\mathrm{~kg}\)”，估计 \(A\) 的概率；
\item 填写下面列联表，并根据列联表判断是否有 \(99\%\) 的把握认为箱产量与养殖方法有关；
\item 根据箱产量的频率分布直方图，对两种养殖方法的优劣进行比较.
\end{enumerate}

\begin{center}
\begin{tabular}{|c|c|c|}
\hline
& 箱产量 \(<50\mathrm{~kg}\) & 箱产量 \(\ge 50\mathrm{~kg}\)\\
\hline
旧养殖法 &  &  \\
\hline
新养殖法 &  &  \\
\hline
\end{tabular}
\end{center}

\begin{center}
\begin{tabular}{|c|c|c|c|}
\hline
\(P(K^2\ge k)\) & 0.050 & 0.010 & 0.001\\
\hline
\(K\) & 3.841 & 6.635 & 10.828\\
\hline
\end{tabular}
\end{center}

附：\(K^2=\dfrac{n(ad-bc)^2}{(a+b)(c+d)(a+c)(b+d)}\).

\texfigure[float=inline,width=0.84\linewidth]{img_repaint/2017/national_paper_2_liberal/q19_fig1.tex}
\end{problem}

\begin{problem}
设 \(O\) 为坐标原点，动点 \(M\) 在椭圆 \(C:\dfrac{x^2}{2}+y^2=1\) 上，过 \(M\) 作 \(x\) 轴的垂线，垂足为 \(N\)，点 \(P\) 满足 \(\overrightarrow{NP}=\sqrt2\,\overrightarrow{NM}\).
\begin{enumerate}
\item 求点 \(P\) 的轨迹方程；
\item 设点 \(Q\) 在直线 \(x=-3\) 上，且 \(\overrightarrow{OP}\cdot\overrightarrow{PQ}=1\). 证明：过点 \(P\) 且垂直于 \(OQ\) 的直线 \(l\) 过 \(C\) 的左焦点 \(F\).
\end{enumerate}
\end{problem}

\begin{problem}
设函数 \(f(x)=(1-x^2)\e^x\).
\begin{enumerate}
\item 讨论 \(f(x)\) 的单调性；
\item 当 \(x\ge0\) 时，\(f(x)\le ax+1\)，求 \(a\) 的取值范围.
\end{enumerate}
\end{problem}

\section{选考题：解答题}

\begin{problem}
在直角坐标系 \(xOy\) 中，以坐标原点为极点，\(x\) 轴的正半轴为极轴建立极坐标系，曲线 \(C_1\) 的极坐标方程为 \(\rho\cos\theta=4\).
\begin{enumerate}
\item \(M\) 为曲线 \(C_1\) 上的动点，点 \(P\) 在线段 \(OM\) 上，且满足 \(\lvert OM\rvert\cdot\lvert OP\rvert=16\)，求点 \(P\) 的轨迹 \(C_2\) 的直角坐标方程；
\item 设点 \(A\) 的极坐标为 \((2,\pi/3)\)，点 \(B\) 在曲线 \(C_2\) 上，求 \(\triangle OAB\) 面积的最大值.
\end{enumerate}
\end{problem}

\begin{problem}
已知 \(a>0\)，\(b>0\)，\(a^3+b^3=2\). 证明：
\begin{enumerate}
\item \((a+b)(a^5+b^5)\ge4\)；
\item \(a+b\le2\).
\end{enumerate}
\end{problem}
